% ============================================================
% SNIPPET PARA CAPÍTULO 4 — Tabla síntesis del protocolo estadístico
% Reemplaza / complementa la descripción textual del protocolo
% en la subsección "Validación estadística y resiliencia operativa"
% ============================================================

\begin{table}[H]
\centering
\caption{Síntesis del protocolo de validación estadística y operativa del Objetivo~4.}
\label{tab:protocolo_obj4}
\footnotesize
\renewcommand{\arraystretch}{1.22}
\begin{tabularx}{\linewidth}{|>{\raggedright\arraybackslash}p{2.8cm}|>{\centering\arraybackslash}p{1.8cm}|>{\centering\arraybackslash}p{2.2cm}|>{\raggedright\arraybackslash}X|}
\hline
\rowcolor{gray!20}
\TableHeaderFirst{Métrica} & \TableHeader{$n$ mínimo} & \TableHeader{Contraste} & \TableHeader{Referencia} \\ \hline

Eficiencia de extracci\'on ($\eta$) &
$n_{\text{pre}}=64$,\par $n_{\text{post}}=38$ &
U de Mann--Whitney unilateral (ejecutado) &
$H_1{:}\;\eta_{\text{post}} > \eta_{\text{pre}}$;\; $\bar{\eta}_{\text{pre}}=4{,}49\%$, $\bar{\eta}_{\text{post}}=22{,}80\%$. \\ \hline

Pérdida de paquetes (\acs{PLR}) &
$\geq 30$ &
U de Mann--Whitney unilateral &
$H_1{:}\;\text{PLR}_{\text{post}} < \text{PLR}_{\text{pre}}$. \\ \hline

\acs{MTTR} broker / red / proceso &
$\geq 30$ por modo &
Descriptiva $+$ \acs{IC} 95\,\% &
Umbral de referencia técnica: $< 60\,\text{s}$. \\ \hline

Latencia extremo a extremo &
$\geq 45$ ciclos &
Descriptiva $+$ \acs{IC} 95\,\% &
Medida desde publicación en gateway hasta recepción en suscriptor. \\ \hline

Soak / estabilidad &
$\geq 6\,\text{h}$ continuas &
Tendencia de memoria y \acs{CPU} &
Umbral: sin degradación acumulada ni spool residual. \\ \hline

Neutralidad \acs{RS-232} &
2 capturas (osciloscopio) &
Cualitativa &
Forma de onda sin perturbación visible en bus original. \\ \hline

\end{tabularx}
\end{table}

La elección del test U de Mann--Whitney como contraste principal se fija
\textit{a priori}, antes de disponer de los datos posteriores. Esta decisión se
justifica por la distribución de la muestra previa: el 54,7\,\% de los valores
de $\eta$ son nulos y la distribución presenta asimetría positiva marcada, lo
que hace improbable superar la normalidad en la prueba de Shapiro--Wilk
($\alpha=0{,}05$). El uso sistemático del contraste no paramétrico elimina la
necesidad de seleccionar retroactivamente el test más favorable.
